%% ODER: format ==         = "\mathrel{==}"
%% ODER: format /=         = "\neq "
%
%
\makeatletter
\@ifundefined{lhs2tex.lhs2tex.sty.read}%
  {\@namedef{lhs2tex.lhs2tex.sty.read}{}%
   \newcommand\SkipToFmtEnd{}%
   \newcommand\EndFmtInput{}%
   \long\def\SkipToFmtEnd#1\EndFmtInput{}%
  }\SkipToFmtEnd

\newcommand\ReadOnlyOnce[1]{\@ifundefined{#1}{\@namedef{#1}{}}\SkipToFmtEnd}
\usepackage{amstext}
\usepackage{amssymb}
\usepackage{stmaryrd}
\DeclareFontFamily{OT1}{cmtex}{}
\DeclareFontShape{OT1}{cmtex}{m}{n}
  {<5><6><7><8>cmtex8
   <9>cmtex9
   <10><10.95><12><14.4><17.28><20.74><24.88>cmtex10}{}
\DeclareFontShape{OT1}{cmtex}{m}{it}
  {<-> ssub * cmtt/m/it}{}
\newcommand{\texfamily}{\fontfamily{cmtex}\selectfont}
\DeclareFontShape{OT1}{cmtt}{bx}{n}
  {<5><6><7><8>cmtt8
   <9>cmbtt9
   <10><10.95><12><14.4><17.28><20.74><24.88>cmbtt10}{}
\DeclareFontShape{OT1}{cmtex}{bx}{n}
  {<-> ssub * cmtt/bx/n}{}
\newcommand{\tex}[1]{\text{\texfamily#1}}	% NEU

\newcommand{\Sp}{\hskip.33334em\relax}


\newcommand{\Conid}[1]{\mathit{#1}}
\newcommand{\Varid}[1]{\mathit{#1}}
\newcommand{\anonymous}{\kern0.06em \vbox{\hrule\@width.5em}}
\newcommand{\plus}{\mathbin{+\!\!\!+}}
\newcommand{\bind}{\mathbin{>\!\!\!>\mkern-6.7mu=}}
\newcommand{\rbind}{\mathbin{=\mkern-6.7mu<\!\!\!<}}% suggested by Neil Mitchell
\newcommand{\sequ}{\mathbin{>\!\!\!>}}
\renewcommand{\leq}{\leqslant}
\renewcommand{\geq}{\geqslant}
\usepackage{polytable}

%mathindent has to be defined
\@ifundefined{mathindent}%
  {\newdimen\mathindent\mathindent\leftmargini}%
  {}%

\def\resethooks{%
  \global\let\SaveRestoreHook\empty
  \global\let\ColumnHook\empty}
\newcommand*{\savecolumns}[1][default]%
  {\g@addto@macro\SaveRestoreHook{\savecolumns[#1]}}
\newcommand*{\restorecolumns}[1][default]%
  {\g@addto@macro\SaveRestoreHook{\restorecolumns[#1]}}
\newcommand*{\aligncolumn}[2]%
  {\g@addto@macro\ColumnHook{\column{#1}{#2}}}

\resethooks

\newcommand{\onelinecommentchars}{\quad-{}- }
\newcommand{\commentbeginchars}{\enskip\{-}
\newcommand{\commentendchars}{-\}\enskip}

\newcommand{\visiblecomments}{%
  \let\onelinecomment=\onelinecommentchars
  \let\commentbegin=\commentbeginchars
  \let\commentend=\commentendchars}

\newcommand{\invisiblecomments}{%
  \let\onelinecomment=\empty
  \let\commentbegin=\empty
  \let\commentend=\empty}

\visiblecomments

\newlength{\blanklineskip}
\setlength{\blanklineskip}{0.66084ex}

\newcommand{\hsindent}[1]{\quad}% default is fixed indentation
\let\hspre\empty
\let\hspost\empty
\newcommand{\NB}{\textbf{NB}}
\newcommand{\Todo}[1]{$\langle$\textbf{To do:}~#1$\rangle$}

\EndFmtInput
\makeatother
%
%
%
%
%
%
% This package provides two environments suitable to take the place
% of hscode, called "plainhscode" and "arrayhscode". 
%
% The plain environment surrounds each code block by vertical space,
% and it uses \abovedisplayskip and \belowdisplayskip to get spacing
% similar to formulas. Note that if these dimensions are changed,
% the spacing around displayed math formulas changes as well.
% All code is indented using \leftskip.
%
% Changed 19.08.2004 to reflect changes in colorcode. Should work with
% CodeGroup.sty.
%
\ReadOnlyOnce{polycode.fmt}%
\makeatletter

\newcommand{\hsnewpar}[1]%
  {{\parskip=0pt\parindent=0pt\par\vskip #1\noindent}}

% can be used, for instance, to redefine the code size, by setting the
% command to \small or something alike
\newcommand{\hscodestyle}{}

% The command \sethscode can be used to switch the code formatting
% behaviour by mapping the hscode environment in the subst directive
% to a new LaTeX environment.

\newcommand{\sethscode}[1]%
  {\expandafter\let\expandafter\hscode\csname #1\endcsname
   \expandafter\let\expandafter\endhscode\csname end#1\endcsname}

% "compatibility" mode restores the non-polycode.fmt layout.

\newenvironment{compathscode}%
  {\par\noindent
   \advance\leftskip\mathindent
   \hscodestyle
   \let\\=\@normalcr
   \let\hspre\(\let\hspost\)%
   \pboxed}%
  {\endpboxed\)%
   \par\noindent
   \ignorespacesafterend}

\newcommand{\compaths}{\sethscode{compathscode}}

% "plain" mode is the proposed default.
% It should now work with \centering.
% This required some changes. The old version
% is still available for reference as oldplainhscode.

\newenvironment{plainhscode}%
  {\hsnewpar\abovedisplayskip
   \advance\leftskip\mathindent
   \hscodestyle
   \let\hspre\(\let\hspost\)%
   \pboxed}%
  {\endpboxed%
   \hsnewpar\belowdisplayskip
   \ignorespacesafterend}

\newenvironment{oldplainhscode}%
  {\hsnewpar\abovedisplayskip
   \advance\leftskip\mathindent
   \hscodestyle
   \let\\=\@normalcr
   \(\pboxed}%
  {\endpboxed\)%
   \hsnewpar\belowdisplayskip
   \ignorespacesafterend}

% Here, we make plainhscode the default environment.

\newcommand{\plainhs}{\sethscode{plainhscode}}
\newcommand{\oldplainhs}{\sethscode{oldplainhscode}}
\plainhs

% The arrayhscode is like plain, but makes use of polytable's
% parray environment which disallows page breaks in code blocks.

\newenvironment{arrayhscode}%
  {\hsnewpar\abovedisplayskip
   \advance\leftskip\mathindent
   \hscodestyle
   \let\\=\@normalcr
   \(\parray}%
  {\endparray\)%
   \hsnewpar\belowdisplayskip
   \ignorespacesafterend}

\newcommand{\arrayhs}{\sethscode{arrayhscode}}

% The mathhscode environment also makes use of polytable's parray 
% environment. It is supposed to be used only inside math mode 
% (I used it to typeset the type rules in my thesis).

\newenvironment{mathhscode}%
  {\parray}{\endparray}

\newcommand{\mathhs}{\sethscode{mathhscode}}

% texths is similar to mathhs, but works in text mode.

\newenvironment{texthscode}%
  {\(\parray}{\endparray\)}

\newcommand{\texths}{\sethscode{texthscode}}

% The framed environment places code in a framed box.

\def\codeframewidth{\arrayrulewidth}
\RequirePackage{calc}

\newenvironment{framedhscode}%
  {\parskip=\abovedisplayskip\par\noindent
   \hscodestyle
   \arrayrulewidth=\codeframewidth
   \tabular{@{}|p{\linewidth-2\arraycolsep-2\arrayrulewidth-2pt}|@{}}%
   \hline\framedhslinecorrect\\{-1.5ex}%
   \let\endoflinesave=\\
   \let\\=\@normalcr
   \(\pboxed}%
  {\endpboxed\)%
   \framedhslinecorrect\endoflinesave{.5ex}\hline
   \endtabular
   \parskip=\belowdisplayskip\par\noindent
   \ignorespacesafterend}

\newcommand{\framedhslinecorrect}[2]%
  {#1[#2]}

\newcommand{\framedhs}{\sethscode{framedhscode}}

% The inlinehscode environment is an experimental environment
% that can be used to typeset displayed code inline.

\newenvironment{inlinehscode}%
  {\(\def\column##1##2{}%
   \let\>\undefined\let\<\undefined\let\\\undefined
   \newcommand\>[1][]{}\newcommand\<[1][]{}\newcommand\\[1][]{}%
   \def\fromto##1##2##3{##3}%
   \def\nextline{}}{\) }%

\newcommand{\inlinehs}{\sethscode{inlinehscode}}

% The joincode environment is a separate environment that
% can be used to surround and thereby connect multiple code
% blocks.

\newenvironment{joincode}%
  {\let\orighscode=\hscode
   \let\origendhscode=\endhscode
   \def\endhscode{\def\hscode{\endgroup\def\@currenvir{hscode}\\}\begingroup}
   %\let\SaveRestoreHook=\empty
   %\let\ColumnHook=\empty
   %\let\resethooks=\empty
   \orighscode\def\hscode{\endgroup\def\@currenvir{hscode}}}%
  {\origendhscode
   \global\let\hscode=\orighscode
   \global\let\endhscode=\origendhscode}%

\makeatother
\EndFmtInput
%
%% ----------------------
%% Stuff for TimCore


%% 'g' inferences




%% --------- Effects ------------


%% ----------------------


%% --------------------------
%%        Desugaring rules
%% --------------------------



%% 'si' is short for "solve or infer" or "unification variable or not" resp

%% 'md' is short for "mode" or "unification variable or not" resp

%% flipsi is a no-op for now



%% ----------- Applications --------------
%% %format applyx a b = a "@" b
%% %format applyx a b = a "\," b

%% Function application with no arguments

%% Function application with 1 argument

%% Function application with 1 argument, no parens

%% Function application with 1 argument, without parenthesis

%% Function application with many arguments

%% ----------- End of Applications --------------



%% Sequences e1; ...; en; v

%% xs etc stole syntax used in examples.

%% Old version: | for BNF, [] for choice
%%%format choice = "[ \! ]"
%%%format `choice` = "\mathop{[ \! ]}"
%% 1mm is about 0.3em with the current font

%% New version: tall | for BNF, fat | for choice
%%format vbar = "\mathop{\bigm|}"
%%format choice = "\boldsymbol{|}"
%%format `choice` = "\mathop{\boldsymbol{|}}"




%% Arrays and tuples
%% %format arrKW = "\textbf{arr}"
%% %format arr (x) = arrKW "\{" x "\}"

%% %format vmap (x) = "\llangle" x "\rrangle"



%% assert = succeed, abbreviated to suc



%%  OLD SYNTAX   %format split (x) (y) (z) = "\textbf{split}(" x ")\{" y "," z "\}"
%% %format defKW = "\textbf{def}"
% only used in comments, but still prettier this way:
%%%%format map = "\textbf{map}"

%% -------Functions ---------------


%% functionOC has open/closed subscript

%% -------End of f unctions ---------------



%% Effects checking, like succeeds{e}


%% Unification
%%%format == = "\!=\!"


% format := = "\mapsto"
%% %format squiggt = "\!\leadsto\!"





































%% ----------- Metatheory --------------

%%%%%%%%%%%%%%%%%%%%%%%%%%%%%%%%%%%%

% ----------------------- Example

\section{Example}
A complete reduction sequence for a small example can be found in \cref{fig:sample-reductions}.
This example shows how constraining the output of a function call can constrain the argument.
While most of the reductions are administrative in nature, these are the highlights:
At \circlednum{1} the \ensuremath{\Varid{swap}} function is inlined so that at \circlednum{2} a \textbeta-reduction can happen.
Step \circlednum{3} inlines the argument, and \circlednum{4} does the matching of the tuple.
At \circlednum{5} and \circlednum{6} the actual numbers are inlined.
\begin{figure}[h]
\begin{footnotesize}
% no space needed around math environments in a figure
\setlength{\abovedisplayskip}{0pt}%
\setlength{\belowdisplayskip}{0pt}%
$$
\begin{array}{@{}l@{\;}l@{\;}l@{}}
                &                    & \ensuremath{\Varid{swap}\langle\Varid{x},\Varid{y}\rangle\hspace{-0.14em}:=\hspace{-0.14em}\langle{}\Varid{y},\Varid{x}\rangle;\,\exists\Varid{p}.\,\Varid{swap}(\Varid{p})\mathrel{=}\langle{}\mathrm{2},\mathrm{3}\rangle;\,\Varid{p}} \\
                & \movestoR{desugar} & \ensuremath{\exists\Varid{swap}.\,\Varid{swap}\hspace{-0.14em}=\hspace{-0.14em}(\lambda\Varid{xy}.\,\exists\Varid{x}\hspace{0.1em}\Varid{y}.\,\langle{}\Varid{x},\Varid{y}\rangle\hspace{-0.14em}=\hspace{-0.14em}\Varid{xy};\,\langle{}\Varid{y},\Varid{x}\rangle);\,\exists\Varid{p}\hspace{0.1em}\Varid{t}.\,\Varid{t}\hspace{-0.14em}=\hspace{-0.14em}\Varid{swap}(\Varid{p});\,\Varid{t}\hspace{-0.14em}=\hspace{-0.14em}\langle{}\mathrm{2},\mathrm{3}\rangle;\,\Varid{p}} \\


\circlednum{1}  & \movestoR{subst,eqn-elim}   & \ensuremath{\exists\Varid{p}\hspace{0.1em}\Varid{t}.\,\Varid{t}\hspace{-0.14em}=\hspace{-0.14em}(\lambda\Varid{xy}.\,\exists\Varid{x}\hspace{0.1em}\Varid{y}.\,\langle{}\Varid{x},\Varid{y}\rangle\hspace{-0.14em}=\hspace{-0.14em}\Varid{xy};\,\langle{}\Varid{y},\Varid{x}\rangle)(\Varid{p});\,\Varid{t}\hspace{-0.14em}=\hspace{-0.14em}\langle{}\mathrm{2},\mathrm{3}\rangle;\,\Varid{p}} \\
                & \movestoR{subst,eqn-elim}   & \ensuremath{\exists\Varid{p}.\,(\mathrm{2},\mathrm{3})\hspace{-0.14em}=\hspace{-0.14em}(\lambda\Varid{xy}.\,\exists\Varid{x}\hspace{0.1em}\Varid{y}.\,\langle{}\Varid{x},\Varid{y}\rangle\hspace{-0.14em}=\hspace{-0.14em}\Varid{xy};\,\langle{}\Varid{y},\Varid{x}\rangle)(\Varid{p});\,\Varid{p}} \\
\circlednum{2}  & \movestoR{app-beta}          & \ensuremath{\exists\Varid{p}.\,(\mathrm{2},\mathrm{3})\hspace{-0.14em}=\hspace{-0.14em}(\exists\Varid{xy}.\,\Varid{xy}\hspace{-0.14em}=\hspace{-0.14em}\Varid{p};\,\exists\Varid{x}\hspace{0.1em}\Varid{y}.\,\langle{}\Varid{x},\Varid{y}\rangle\hspace{-0.14em}=\hspace{-0.14em}\Varid{xy};\,\langle{}\Varid{y},\Varid{x}\rangle);\,\Varid{p}} \\
                & \movestoR{exi-float}          & \ensuremath{\exists\Varid{p}\hspace{0.1em}\Varid{xy}.\,(\mathrm{2},\mathrm{3})\hspace{-0.14em}=\hspace{-0.14em}((\Varid{xy}\hspace{-0.14em}=\hspace{-0.14em}\Varid{p};\,\exists\Varid{x}\hspace{0.1em}\Varid{y}.\,\langle{}\Varid{x},\Varid{y}\rangle\hspace{-0.14em}=\hspace{-0.14em}\Varid{xy};\,\langle{}\Varid{y},\Varid{x}\rangle));\,\Varid{p}} \\
\circlednum{3}  & \movestoR{subst,eqn-elim}   & \ensuremath{\exists\Varid{p}.\,(\mathrm{2},\mathrm{3})\hspace{-0.14em}=\hspace{-0.14em}(\exists\Varid{x}\hspace{0.1em}\Varid{y}.\,\langle{}\Varid{x},\Varid{y}\rangle\hspace{-0.14em}=\hspace{-0.14em}\Varid{p};\,\langle{}\Varid{y},\Varid{x}\rangle);\,\Varid{p}} \\
                & \movestoR{exi-float,exi-float} & \ensuremath{\exists\Varid{p}\hspace{0.1em}\Varid{x}\hspace{0.1em}\Varid{y}.\,(\mathrm{2},\mathrm{3})\hspace{-0.14em}=\hspace{-0.14em}(\langle{}\Varid{x},\Varid{y}\rangle\hspace{-0.14em}=\hspace{-0.14em}\Varid{p};\,\langle{}\Varid{y},\Varid{x}\rangle);\,\Varid{p}} \\
                & \movestoR{eqn-float,seq-assoc}& \ensuremath{\exists\Varid{p}\hspace{0.1em}\Varid{x}\hspace{0.1em}\Varid{y}.\,\langle{}\Varid{x},\Varid{y}\rangle\hspace{-0.14em}=\hspace{-0.14em}\Varid{p};\,(\mathrm{2},\mathrm{3})\hspace{-0.14em}=\hspace{-0.14em}\langle{}\Varid{y},\Varid{x}\rangle;\,\Varid{p}} \\
                & \movestoR{hnf-swap}          & \ensuremath{\exists\Varid{p}\hspace{0.1em}\Varid{x}\hspace{0.1em}\Varid{y}.\,\Varid{p}\hspace{-0.14em}=\hspace{-0.14em}\langle{}\Varid{x},\Varid{y}\rangle;\,(\mathrm{2},\mathrm{3})\hspace{-0.14em}=\hspace{-0.14em}\langle{}\Varid{y},\Varid{x}\rangle;\,\Varid{p}} \\
                & \movestoR{subst,eqn-elim}   & \ensuremath{\exists\Varid{x}\hspace{0.1em}\Varid{y}.\,(\mathrm{2},\mathrm{3})\hspace{-0.14em}=\hspace{-0.14em}\langle{}\Varid{y},\Varid{x}\rangle;\,\langle{}\Varid{x},\Varid{y}\rangle} \\
\circlednum{4}  & \movestoR{u-tup,seq-assoc}    & \ensuremath{\exists\Varid{x}\hspace{0.1em}\Varid{y}.\,\mathrm{2}\hspace{-0.14em}=\hspace{-0.14em}\Varid{y};\,\mathrm{3}\hspace{-0.14em}=\hspace{-0.14em}\Varid{x};\,\langle{}\Varid{x},\Varid{y}\rangle} \\
                & \movestoR{hnf-swap}          & \ensuremath{\exists\Varid{x}\hspace{0.1em}\Varid{y}.\,\Varid{y}\hspace{-0.14em}=\hspace{-0.14em}\mathrm{2};\,\mathrm{3}\hspace{-0.14em}=\hspace{-0.14em}\Varid{x};\,\langle{}\Varid{x},\Varid{y}\rangle} \\
\circlednum{5}  & \movestoR{subst,eqn-elim}   & \ensuremath{\exists\Varid{x}.\,\mathrm{3}\hspace{-0.14em}=\hspace{-0.14em}\Varid{x};\,\langle{}\Varid{x},\mathrm{2}\rangle} \\
                & \movestoR{hnf-swap}          & \ensuremath{\exists\Varid{x}.\,\Varid{x}\hspace{-0.14em}=\hspace{-0.14em}\mathrm{3};\,\langle{}\Varid{x},\mathrm{2}\rangle} \\
\circlednum{6}  & \movestoR{subst,eqn-elim}   & \ensuremath{\langle{}\mathrm{3},\mathrm{2}\rangle} \\
\end{array}
$$
\end{footnotesize}
\caption{A sample reduction sequence} \label{fig:sample-reductions}
\Description{A sample reduction sequence}
\end{figure}

